% Imporditud: tools/import_eksperimendid.py
% Allikas: Eesti füüsikaolümpiaadi eksperimendi ülesannete
%   kogu (Rael Kalda, 2023), ülesanne Ü172, lahendus L172.
\setAuthor{Jaan Kalda}
\setRound{lõppvoor}
\setYear{2016}
\setNumber{G E2}
\setDifficulty{5}
\setTopic{dünaamika}
\setVahendid{Purk veega, õhukeseseinaline kilekott, väike raskus massiga m = 8,5 g (tihedusega $\rho$k = 7,8 g/cm$^3$), ühtlase laiusega paberiribad, joonlaud, kuivatuspaber}
\setPunktid{14}
\setAllikas{Eksperimendi kogu Ü172, lahendus lk 130}
\setLahendusseis{toores}

\prob{Kile ja paber}
Leidke kile ja paberi vaheline hõõrdetegur.

\hint

\solu
Paneme tina kilekotti ja paneme ka pabeririba otsapidi kotti ning sukeldame sel-

le vette. Leiame, kui pikalt (H) peab olema pabeririba veenivoost allpool, et tina

koos kotiga kerkiks üles, kui pabeririba otsast tirida. Mõõdame riba laiuse L. Keskmine rõhumisjõud pabeririba ühele küljele on $\rho$vgH2L/2 ning vastav hõõrdejõud --- $\mu\rho$vgH2L/2. Et ribal on kaks külge, siis summaarne jõud on kaks korda suurem, $\mu\rho$vgH2L, mis peab olema võrdne haavlile mõjuva raskusjõu ja üleslükkejõu

vahega, $\mu\rho$vgH2L = mg(1 $- \rho$v/$\rho$t),

millest saame m($\rho-$v 1 $- \rho$t$-$1) H 2L $\mu$ = .
\probend
